\documentclass[12pt,a4paper]{ctexart}

\usepackage{amsmath,amssymb}
\usepackage{geometry}
\geometry{left=2.2cm,right=2.2cm,top=2.2cm,bottom=2.2cm}

\setlength{\parindent}{0pt}
\setlength{\parskip}{0.6em}

\title{大学物理（2）期末线上 2022 春\\计算题答案与解析}
\author{}
\date{}

\begin{document}
\maketitle

\section*{第 1 题：椭圆偏振光通过偏振片}

\subsection*{题目模型}

设椭圆偏振光由 $x$ 方向和 $y$ 方向两束相干线偏振光合成。题目给出 $y$ 方向振幅与 $x$ 方向振幅之比为 $2:1$，因此可设
\[
E_x=A\cos\omega t,\qquad
E_y=2A\cos(\omega t+\varphi).
\]

偏振片通光方向与 $x$ 轴夹角为 $\theta$，所以通过偏振片后的电场是原电场在通光方向上的投影：
\[
E=E_x\cos\theta+E_y\sin\theta.
\]

\subsection*{光强计算}

光强正比于电场平方的时间平均：
\[
I\propto \langle E^2\rangle.
\]

展开：
\[
\begin{aligned}
\langle E^2\rangle
&=\langle E_x^2\rangle\cos^2\theta
 +\langle E_y^2\rangle\sin^2\theta
 +2\langle E_xE_y\rangle\sin\theta\cos\theta.
\end{aligned}
\]

分别计算三个平均值：
\[
\langle E_x^2\rangle=\frac{A^2}{2},
\]
\[
\langle E_y^2\rangle=\frac{(2A)^2}{2}=2A^2,
\]
\[
\langle E_xE_y\rangle
=2A^2\langle \cos\omega t\cos(\omega t+\varphi)\rangle
=A^2\cos\varphi.
\]

因此
\[
\langle E^2\rangle
=\frac{A^2}{2}\cos^2\theta
+2A^2\sin^2\theta
+2A^2\cos\varphi\sin\theta\cos\theta.
\]

原椭圆偏振光总光强为 $I_0$，而
\[
I_0\propto \langle E_x^2\rangle+\langle E_y^2\rangle
=\frac{A^2}{2}+2A^2
=\frac{5A^2}{2}.
\]

所以通过偏振片后的光强为
\[
\boxed{
I=\frac{I_0}{5}
\left(
\cos^2\theta+4\sin^2\theta+4\sin\theta\cos\theta\cos\varphi
\right)
}
\]

也可写成
\[
\boxed{
I=\frac{I_0}{5}
\left(
\cos^2\theta+4\sin^2\theta+2\sin2\theta\cos\varphi
\right)
}.
\]

\subsection*{得分点}

\begin{itemize}
  \item 写出 $E_x=A\cos\omega t,\ E_y=2A\cos(\omega t+\varphi)$。
  \item 写出偏振片投影 $E=E_x\cos\theta+E_y\sin\theta$。
  \item 光强必须取时间平均 $\langle E^2\rangle$。
  \item 最后要用总光强 $I_0\propto 5A^2/2$ 归一化。
\end{itemize}

\section*{第 2 题：动量本征态展开}

\subsection*{题目模型}

题给波函数
\[
\Psi(x)=\frac{4\cos^2kx+\sqrt{2}-2}{2\sqrt{2\pi\hbar}},
\]
动量本征函数为
\[
\Phi_{p_x}(x)=\frac{1}{\sqrt{2\pi\hbar}}e^{\frac{i}{\hbar}p_xx}.
\]

要求测量动量时可能得到的取值、对应概率，以及动量平均值。

\subsection*{化简波函数}

利用恒等式
\[
\cos^2kx=\frac{1+\cos2kx}{2}.
\]

于是
\[
\begin{aligned}
4\cos^2kx+\sqrt{2}-2
&=4\cdot\frac{1+\cos2kx}{2}+\sqrt{2}-2\\
&=2+2\cos2kx+\sqrt{2}-2\\
&=\sqrt{2}+2\cos2kx.
\end{aligned}
\]

因此
\[
\Psi(x)=\frac{\sqrt{2}+2\cos2kx}{2\sqrt{2\pi\hbar}}
=\frac{1}{\sqrt{2}}\frac{1}{\sqrt{2\pi\hbar}}
+\frac{\cos2kx}{\sqrt{2\pi\hbar}}.
\]

再用
\[
\cos2kx=\frac{e^{i2kx}+e^{-i2kx}}{2}.
\]

所以
\[
\Psi(x)
=\frac{1}{\sqrt{2}}\frac{1}{\sqrt{2\pi\hbar}}
+\frac{1}{2}\frac{e^{i2kx}}{\sqrt{2\pi\hbar}}
+\frac{1}{2}\frac{e^{-i2kx}}{\sqrt{2\pi\hbar}}.
\]

和动量本征函数比较：
\[
\frac{1}{\sqrt{2\pi\hbar}}=\Phi_0(x),
\]
\[
\frac{e^{i2kx}}{\sqrt{2\pi\hbar}}=\Phi_{2\hbar k}(x),
\]
\[
\frac{e^{-i2kx}}{\sqrt{2\pi\hbar}}=\Phi_{-2\hbar k}(x).
\]

因此
\[
\boxed{
\Psi(x)=
\frac{1}{\sqrt{2}}\Phi_0(x)
+\frac{1}{2}\Phi_{2\hbar k}(x)
+\frac{1}{2}\Phi_{-2\hbar k}(x)
}.
\]

\subsection*{动量可能值和概率}

展开系数的模平方就是测得对应动量的概率：
\[
P(p_x=0)=\left|\frac{1}{\sqrt{2}}\right|^2=\frac{1}{2},
\]
\[
P(p_x=2\hbar k)=\left|\frac{1}{2}\right|^2=\frac{1}{4},
\]
\[
P(p_x=-2\hbar k)=\left|\frac{1}{2}\right|^2=\frac{1}{4}.
\]

所以可能动量为
\[
\boxed{
p_x=0,\quad 2\hbar k,\quad -2\hbar k
}
\]
对应概率为
\[
\boxed{
\frac{1}{2},\quad \frac{1}{4},\quad \frac{1}{4}
}.
\]

\subsection*{动量平均值}

\[
\begin{aligned}
\langle p_x\rangle
&=0\cdot\frac{1}{2}
+2\hbar k\cdot\frac{1}{4}
+(-2\hbar k)\cdot\frac{1}{4}\\
&=0.
\end{aligned}
\]

因此
\[
\boxed{\langle p_x\rangle=0}.
\]

\subsection*{得分点}

\begin{itemize}
  \item 先把 $\cos^2kx$ 化成 $\frac{1+\cos2kx}{2}$。
  \item 再把 $\cos2kx$ 化成指数形式。
  \item 和 $e^{ip_xx/\hbar}$ 对比，读出动量。
  \item 概率是展开系数模平方，不是系数本身。
\end{itemize}

\section*{第 3 题：$N$ 缝光栅夫琅禾费衍射光强公式}

\subsection*{定义}

设单缝宽度为 $a$，光栅常数为 $d$，正入射，衍射角为 $\theta$。

定义
\[
\beta=\frac{\pi a\sin\theta}{\lambda},
\qquad
\alpha=\frac{\pi d\sin\theta}{\lambda}.
\]

其中：
\begin{itemize}
  \item $\beta$ 由单缝宽度 $a$ 决定，控制单缝衍射包络；
  \item $\alpha$ 由光栅常数 $d$ 决定，控制多缝干涉因子。
\end{itemize}

\subsection*{（1）单缝衍射振幅公式}

题目给单缝中央明纹中心处振幅为 $E_0$，强度为 $I_0$。

单缝夫琅禾费衍射在衍射角 $\theta$ 方向的振幅为
\[
\boxed{
E_1(\theta)=E_0\frac{\sin\beta}{\beta}
}.
\]

因此单缝光强为
\[
\boxed{
I_1(\theta)=I_0\left(\frac{\sin\beta}{\beta}\right)^2
}.
\]

\subsection*{（2）$N$ 缝光栅衍射光强公式}

相邻两缝在方向 $\theta$ 上的光程差为
\[
\Delta=d\sin\theta.
\]

对应相位差为
\[
\delta=\frac{2\pi}{\lambda}d\sin\theta=2\alpha.
\]

$N$ 个等振幅、等相位差的波叠加，其振幅因子为
\[
\frac{\sin N\alpha}{\sin\alpha}.
\]

因此 $N$ 缝光栅的总振幅为
\[
\boxed{
E(\theta)=
E_0\frac{\sin\beta}{\beta}
\frac{\sin N\alpha}{\sin\alpha}
}.
\]

光强与振幅平方成正比，所以
\[
\boxed{
I(\theta)=
I_0
\left(\frac{\sin\beta}{\beta}\right)^2
\left(\frac{\sin N\alpha}{\sin\alpha}\right)^2
}.
\]

\subsection*{补充结论}

主极大条件：
\[
d\sin\theta=k\lambda.
\]

单缝暗纹条件：
\[
a\sin\theta=m\lambda.
\]

缺级条件：
\[
d\sin\theta=k\lambda
\quad\text{且}\quad
a\sin\theta=m\lambda
\]
同时成立。

\subsection*{得分点}

\begin{itemize}
  \item 写清楚 $\beta=\frac{\pi a\sin\theta}{\lambda}$，$\alpha=\frac{\pi d\sin\theta}{\lambda}$。
  \item 单缝因子是 $\frac{\sin\beta}{\beta}$。
  \item 多缝干涉因子是 $\frac{\sin N\alpha}{\sin\alpha}$。
  \item 光强公式中两个因子都要平方。
\end{itemize}

\end{document}
